\chapter{El cuerpo $\displaystyle \C $}
\section{Definiciones básicas}
Informalmente, podemos definir el conjunto $\displaystyle \C $ como
\[ \C = \left\{ a + bi \: : \; a,b \in \R\right\}  ,\]
tal que $\displaystyle i^{2}=-1 $. Construyamos $\displaystyle \C $ de una manera más satisfactoria.
\begin{definition}[El conjunto $\displaystyle \C $]
Definimos $\displaystyle \C $ como el conjunto
\[ \C : = \R[\mathtt{x}] / \left(\mathtt{x}^{2} + 1\right).\]
Asímismo, definimos la suma y el producto en $\displaystyle \C $ están definidos como lo están en un cociente de anillos.
\end{definition}
Alternativamente, podemos definir $\displaystyle \C $ como un $\displaystyle \R $-espacio vectorial de dimensión 2, con la suma usual de vectores. Añadimos además la operación del \textbf{producto complejo:}
\[\left(a,b\right) \cdot \left(c,d\right) := \left(ac-bd, ad + bc\right) .\]
De esta forma tenemos que $\displaystyle \left(0,1\right) \cdot \left(0,1\right)=\left(-1, 0\right) $. Esta notación es engorrosa, por lo que escribiremos $\displaystyle \left(a,b\right)= a + bi $. 

\begin{prop}
	El trío $\displaystyle \left(\C, +, \cdot \right) $ es un cuerpo conmutativo, del cual $\displaystyle \R $ es un subcuerpo y está identificado por $\displaystyle \left\{ x + iy \in \C \; : \; y = 0\right\}  $.
\end{prop}
\begin{proof}
Sean $\displaystyle z,w,\xi \in \C $.
\begin{enumerate}
\item Asociatividad de la suma. Aprovechando la asociatividad de la suma en $\displaystyle \R $,
	\[
	\begin{split}
		z + \left(w + \xi \right) = & \left(a + bi\right) + \left(\left(c + di\right)+\left(e + fi\right)\right) = \left(a + bi\right)+\left(\left(c + e\right)+\left(d + f\right)i\right) \\
		= & \left(a + c + e\right) + \left(b + d + f\right)i= \left(\left(a + c\right)+\left(b + d\right)i\right)+\left(e + fi\right)\\
		= &  \left(z + w\right) + \xi .
	\end{split}
	\]
\item Conmutatividad de la suma. Aprovechando la conmutatividad de la suma en $\displaystyle \R $,
	\[
	\begin{split}
		z + w = & \left(a + bi\right)+\left(c + di\right)=\left(a + c\right) + \left(b + d\right)i = \left(c + a\right) + \left(d + b\right)i \\
		= & \left(c + di\right) + \left(a +bi\right)=w + z  .
	\end{split}
	\]
\item El $\displaystyle 0 $ es el elemento neutro para la suma:
	\[z + 0 = \left(a + bi\right) + \left(0 + 0i\right)= \left(a + 0\right)+\left(b+0\right)i = a + bi = z.\]
\item Existencia de opuestos para la suma. Si $\displaystyle z = a + bi $, basta con tomar $\displaystyle -z = -a-bi $, en efecto
	\[z + \left(-z\right)= \left(a + bi\right)+\left(-a-bi\right)=\left(a-a\right)+\left(b-b\right)i=0 .\]
\item Asociatividad del producto.
	\[
	\begin{split}
		z\left(w\xi\right) = & \left(a + bi\right)\left(\left(c+di\right)\left(e + fi\right)\right)=\left(a + bi\right)\left(\left(ce-df\right)+\left(cd+ed\right)i\right)\\
		= & \left(ace - adf -bcd - bed\right) + \left(bce - bdf + acd + aed\right)i \\
		= & \left(\left(ac-db\right)+\left(ad + bd\right)i\right)\left(e + fi\right) = \left(zw\right)\xi .
	\end{split}
	\]
\item Conmutatividad del producto.
	\[zw = \left(a + bi\right)\left(c + di\right)=\left(ac - bd\right) + \left(ad + bc\right)i = \left(ca - db\right)+\left(da + cb\right)i = wz .\]
\item El $\displaystyle 1 $ es el elemento neutro para el producto.
	\[1z = \left(1 + 0i\right)\left(a + bi\right)=\left(a - 0\right) + \left(b + 0\right)i = a + bi = z .\]
\item Existencia de inverso para el producto. Si $\displaystyle z = a + bi \neq 0 $, entonces es sencillo comprobar que si tomamos $\displaystyle z^{-1}=\frac{a -bi}{a^{2}+b^{2}} $ se cumple que $\displaystyle zz^{-1}= 1 $.
\item Propiedad distributiva del producto respecto de la suma. La demostración es engorrosa por lo que no se hará aquí.
\end{enumerate}
\end{proof}
 
\begin{definition}
Dado $\displaystyle z \in \C $, llamamos a $\displaystyle \overline{z} $ su \textbf{complejo conjugado} de forma que si $\displaystyle z = a + bi $, entonces $\displaystyle \overline{z}= a-bi $. 
\end{definition}
\begin{observation}
La conjugación $\displaystyle z \to \overline{z} $ es una reflexión de $\displaystyle \C $ respecto de la recta $\displaystyle \R $. 
\end{observation}
\begin{definition}[Módulo]
El \textbf{módulo} de $\displaystyle z = a + bi \in \C $ viene dado por
\[ \left|z\right|=\sqrt{a^{2}+b^{2}} .\]
\end{definition}

\begin{notation}[Plano complejo]
	Dado un número complejo $\displaystyle z = a + bi $, decimos que su \textbf{parte real} es $\displaystyle a = \Rel z $ y su \textbf{parte imaginaria} es $\displaystyle b = \Imagen z $. Geométricamente, tenemos que $\displaystyle \Rel z $ es la proyección ortogonal de $\displaystyle z $ sobre la recta $\displaystyle \R $ en $\displaystyle \C $, y $\displaystyle i \Imagen z $ es la proyección ortogonal de $\displaystyle z $ sobre la recta $\displaystyle i\R := \left\{ x + iy \: : \; x = 0\right\}  $.
\end{notation}

\begin{prop}
	Sean $\displaystyle z,w \in \C $.
\begin{enumerate}
\item $\displaystyle z^{-1}=\frac{\overline{z}}{ \left|z\right|^{2}} $.
\item $\displaystyle \Rel z =\frac{z + \overline{z}}{2} $. 
\item $\displaystyle \Imagen z=\frac{z-\overline{z}}{2i} $. 
\item $\displaystyle \overline{z + w}= \overline{z} + \overline{w} $.
\item $\displaystyle \overline{zw} = \overline{z} \cdot \overline{w} $.
\item $\displaystyle \left|z\right|= \left|\overline{z}\right| $.
\item $\displaystyle \left|z\right|^{2} = z \cdot \overline{z} $.
\item $\displaystyle \left|zw\right| = \left|z\right| \left|w\right| $.
\end{enumerate}
\end{prop}
\begin{proof}
Sean $\displaystyle z = a +bi, w = c + di \in \C $. 
\begin{enumerate}
\item Es trivial a partir de la definición de inversa que dimos en la demostración anterior.
\item Si $\displaystyle z = a + bi $ tenemos que 
	\[\frac{z + \overline{z}}{2}=\frac{\left(a + bi\right)+\left(a-bi\right)}{2}=\frac{2a}{2}=a = \Rel z .\]
\item Si $\displaystyle z = a + bi $ tenemos que 
	\[\frac{z -\overline{z}}{2i}=\frac{\left(a+bi\right)-\left(a-bi\right)}{2i}=\frac{2bi}{2i}=b=\Imagen z .\]
\item Sean $\displaystyle z = a + bi $ y $\displaystyle w = c + di $, entonces
	\[\overline{z + w}= \left(a + c\right)-\left(b + d\right)i = \left(a - bi\right)+\left(c-di\right) = \overline{z} + \overline{w} .\]
\item Si $\displaystyle z = a + bi $ y $\displaystyle w = c + di $, entonces
	\[\overline{zw}=\left(ac - bd\right)-\left(ad + bc\right)i .\]
	Por otro lado, 
	\[\overline{z} \cdot \overline{w}= \left(a -bi\right)\left(c -di\right)=\left(ac-bd\right)+\left(-ad-bc\right)i .\]
\item Supongamos que $\displaystyle z = a + bi $, 
	\[ \left|\overline{z}\right|=\sqrt{a^{2}+\left(-b\right)^{2}}=\sqrt{a^{2}+b^{2}}= \left|z\right| .\]
\item Si $\displaystyle z = a + bi $,
	\[ z \cdot \overline{z}=\left(a + bi\right)\left(a - bi\right)=\left(a^{2}+b^{2}\right)+\left(-ab + ab\right)i = \left(\sqrt{a^{2}+b^{2}}\right)^{2}= \left|z\right|^{2} .\]
\item Sean $\displaystyle z = a + bi $ y $\displaystyle w = c + di $,
	\[ \left|zw\right|=\sqrt{\left(ac - bd\right)^{2}+\left(ad + bc\right)^{2}} = \sqrt{a^{2}c^{2} + b^{2}d^{2}+a^{2}d^{2}+b^{2}c^{2}}.\]
Por otro lado, 
\[ \left|z\right| \left|w\right|=\sqrt{a^{2}+b^{2}}\sqrt{c^{2}+d^{2}}=\sqrt{a^{2}c^{2}+a^{2}d^{2}+b^{2}c^{2}+b^{2}d^{2}} .\]
\end{enumerate}
\end{proof}
\begin{notation}[Coordenadas polares]
Las coordenadas polares explotan simetrías radiales, por lo que son muy útiles. Cada número complejo $\displaystyle z = x + yi $ se puede expresar de la forma 
\[z = x + iy = r\left(\cos\theta + i \sin \theta \right) ,\]
con $\displaystyle r = \sqrt{x^{2}+y^{2}} $ y $\displaystyle \theta = \arctan\frac{y}{x} $. Podmeos ver que $\displaystyle \theta $ está determinado salvo suma de un múltiplo entero de $\displaystyle 2\pi  $. Para que haya biyección entre ambas representaciones nos restringimos de la siguiente manera
\[\left(0,\infty\right)\times\left(-\pi, \pi \right] \to \C/ \left\{ 0\right\} : \left(r, \theta \right)\to r\left(\cos\theta + i \sin \theta\right) .\]
\end{notation}

\section{Funciones elementales}
\begin{definition}[Exponencial compleja]
Dado $\displaystyle \theta \in \R $, se define
\[e^{i\theta }:= \cos \theta + i \sin \theta  .\]
Más generalmente, si $\displaystyle z = x + iy $ ponemos
\[e^{z}:= e^{x +iy}=e^{x} \cdot e^{iy} .\]
\end{definition}

Así, podemos escribir la forma polar de un número complejo como
\[z = re^{i\theta} ,\]
siendo $\displaystyle r = \left|z\right| $ y $\displaystyle \theta $ está determinado unívocamente en $\displaystyle \theta \in \left(-\pi , \pi \right]  $ para $\displaystyle z \neq 0 $. 
\begin{notation}
Diremos que $\displaystyle \theta = \Arg\left(z \right) $. Diremos que $\displaystyle \arg\left(z\right)= \left\{ \Arg\left(z\right)+2k\pi \; : \; k \in \Z\right\}  $. A la aplicación $\displaystyle z \to \Arg z $ la llamaremos \textbf{rama principal del argumento}. 
\end{notation}
\begin{definition}[Logaritmo complejo]
Definimos la \textbf{rama principal del logaritmo} como
\[\log : \C / (-\infty, 0] \to \C : z \to \log \left|z\right| + i \Arg\left(z\right) .\]
\end{definition}
\begin{observation}
	Podemos hacer las siguientes observaciones.
	\begin{enumerate}
		\item Se tiene que $\displaystyle e ^{\log z}=z $ y que $\displaystyle \log z $ coincide con el logaritmo usual en $\displaystyle \R $ cuando $\displaystyle z \in \left(0, \infty\right) $.
		\item La función logaritmo definida de esta manera es inyectiva, pero la función exponencial no lo es. En efecto, $\displaystyle e^{z + 2k\pi i}=e^{z} $, $\displaystyle \forall k \in \Z $. 
	\end{enumerate} 
\end{observation}

\begin{prop}
Sean $\displaystyle \theta, \varphi \in \R $ y $\displaystyle z,w \in \C $. 
\begin{enumerate}
\item $\displaystyle \left|e^{i\theta }\right|=1 $.
\item $\displaystyle \overline{e^{i\theta }}=e^{-i\theta }=\frac{1}{e^{i\theta}} $. 
\item $\displaystyle e^{i\left(\theta + \varphi\right)}=e^{i\theta } \cdot e^{i\varphi} $.
\item $\displaystyle e^{z + w}=e^{z} \cdot e^{w} $. 
\item $\displaystyle \arg \overline{z} = - \arg z $ para $\displaystyle z \neq 0 $. 
\item Si $\displaystyle z \neq 0 $, $\displaystyle \arg\left(\frac{1}{z}\right)=-\arg z $.
\item $\displaystyle \arg\left(z \cdot w\right)=\arg z + \arg w $. 
\end{enumerate}
\end{prop}
\begin{proof}
Supongamos que $\displaystyle \theta, \varphi \in \R $ y $\displaystyle z, w \in \C $.
\begin{enumerate}
\item Tenemos que $\displaystyle \left|e^{i\theta }\right|=1 $ es equivalente a que $\displaystyle \cos^{2}\theta + \sin ^{2}\theta = 1 $. 
\item Es sencillo ver que 
	\[\overline{e^{i\theta}} = \cos\theta - i \sin \theta = \cos \theta + i \sin \left(-\theta \right)=e^{-i\theta} .\]
	Además, como $\displaystyle z^{-1}= \frac{\overline{z}}{ \left|z\right|^{2}} $, tenemos que $\displaystyle e^{-i\theta}=\frac{1}{e^{i\theta}} $. 
\item Aplicando las fórmulas trigonométricas
	\[
	\begin{split}
		e^{i\left(\theta + \varphi\right)} = & \cos\left(\theta + \varphi\right) + i \sin \left(\theta + \varphi\right) = \cos\theta \cos\varphi - \sin \theta \sin \varphi + i \left(\sin\theta \cos\varphi + \sin\varphi \cos \theta\right) \\
		= & \cos\theta\left(\cos\varphi + i\sin\varphi\right)+i\sin\theta\left(\cos\varphi + i \sin \varphi\right) = \left(\cos\theta + i \sin \theta\right)\left(\cos\varphi + i\sin\varphi\right)\\
		= & e^{i\theta}e^{i\varphi} .
	\end{split}
	\]
\item Si escribimos $\displaystyle z = a + bi $ y $\displaystyle w = x + iy $ con $\displaystyle a,b,x,y \in \R $, tendremos que aplicando \textbf{(3)} 
	\[e^{z + w}= e^{\left(a + x\right) + \left(b + y\right)i} = e^{a + x} \cdot e^{\left(b + y\right)i} = e^{a}e^{x} e^{bi}e^{yi}=e^{a + bi} \cdot e^{x + yi}= e^{z} \cdot e^{w}.\]
\item A partir de \textbf{(2)} es sencillo ver que si $\displaystyle z = re^{i\theta} $, 
	\[\overline{z}=\overline{re^{i\theta}}=r\overline{e^{i\theta}}=re^{-i\theta} .\]
\item Es trivial a partir de \textbf{(2)}.
\item Si $\displaystyle z = \left|z\right|e^{i\theta} $ y $\displaystyle w = \left|w\right|e^{i\varphi} $, por \textbf{(4)} tenemos que 
	\[z \cdot w = \left|z\right| \left|w\right|e^{i\theta} \cdot e^{i\varphi}= \left|z\right| \left|w\right| e^{i\left(\theta + \varphi\right)} .\]
\end{enumerate}
\end{proof}

\begin{observation}[Interpretación geométrica del producto complejo]
Dado $\displaystyle w \in \C $, consideremos la función $\displaystyle z \to w \cdot z $. Utilizando la forma polar, tenemos que 
\[w \cdot z = \left(s e^{i\varphi}\right)\left(re^{i\theta }\right)=sre^{i\left(\varphi+\theta \right)} .\]
Como el producto es conmutativo, tenemos que la composición de homotecias y rotaciones también es conmutativa.
\end{observation}
Estos resultados nos permiten calcular las raíces $\displaystyle n $-ésimas de un número complejo $\displaystyle w $ para cualquier $\displaystyle n \in \N $. En efecto, si $\displaystyle w = \rho e^{i\varphi} $ y $\displaystyle z = re^{i\theta} $, planteamos la ecuación $\displaystyle z^{n}= w $ en forma polar,
\[r^{n}e^{in\theta}=\rho e^{i\varphi} .\]
Así, igualando módulos y argumentos deducimos que $\displaystyle r = \rho^{\frac{1}{n}} $ y $\displaystyle \varphi = n\theta + 2k\pi  $ con $\displaystyle k \in \Z $. De esta forma, podemos afirmar que las $\displaystyle n $ raíces $\displaystyle n $-ésimas de $\displaystyle w $ son los números $\displaystyle z_{j} $ que tienen la forma
\[z_{j}=\rho^{\frac{1}{n}}e^{i\theta_{j}}, \; \theta_{j}=\frac{\varphi + 2j\pi }{n}, j = 0, 1, \ldots, n-1 .\]

\begin{definition}[Raíz $\displaystyle n $-ésima compleja]
Vamos a definir
\[f_{0} : \C / (-\infty, 0] \to \C : w \to e^{\frac{1}{n}\log w} .\]
Donde $\displaystyle \log w $ es la rama principal del logaritmo complejo.
\end{definition}
Las otras raíces se obtienen según
\[f_{j}\left(w\right)=f_{0}\left(w\right)e^{i\frac{j2\pi }{n}}, \; j = 1, \ldots, n-1 .\]
Para todas estas funciones con $\displaystyle j = 0, \ldots, n-1 $ 
\[f_{j} : \C / (-\infty, 0] \to U_{j} = \left\{ z \in \C \; : \; arg z \in \left(-\frac{\pi }{n}, \frac{\pi }{n}\right)+\frac{2\pi j}{n}\right\}  .\]
La función $\displaystyle z^{n} $ es inyectiva y $\displaystyle f_{j} $ es su inversa. 
La función de $\displaystyle f_{0} $ inspira la siguiente definición de $\displaystyle z^{\alpha } $ potencia compleja. Concretamente, dado $\displaystyle \alpha \in \C $, 
\[\varphi_{\alpha }: \C/ (-\infty, 0] \to \C : z \to z^{\alpha }:=e^{\alpha\log z} ,\]
es la potencia compleja $\displaystyle \alpha  $. 
\begin{definition}[Funciones trigonométricas]
Vamos a definir
\[\cos z := \frac{e^{iz}+e^{-iz}}{2} \quad \text{y} \quad \sin z : = \frac{e^{iz}-e^{-iz}}{2i} .\]
\end{definition}
\begin{observation}
Para preservar diferenciabilidad (en sentido complejo) es la única manera de definir las funciones trigonométricas.
\end{observation}
\begin{observation}
Las funciones trigonométricas complejas no son acotadas en $\displaystyle \C $.
\end{observation}
\begin{prop}
Sean $\displaystyle k \in \Z $ y $\displaystyle z \in \C $.
\begin{enumerate}
\item $\displaystyle \cos\left(-z\right)= \cos z $.
\item $\displaystyle \sin\left(-z\right)=-\sin z $.
\item $\displaystyle \sin \left(z + 2k\pi \right)= \sin z $.
\item $\displaystyle \cos\left(z + 2k\pi \right)= \cos z $.
\item $\displaystyle \cos z = \sin\left(z + \frac{\pi }{2}\right) $. 
\end{enumerate}
\end{prop}
\begin{notation}
Las funciones de tangente, cotangente, etc, se definen a partir de la del seno y coseno como lo hacen en el caso real.
\end{notation}
\begin{definition}[Funciones hiperbólicas]
Vamos a definir 
\[\cosh z := \frac{e^{z}+e^{-z}}{2} \quad \text{y} \quad \sinh z := \frac{e^{z}-e^{-z}}{2} .\]
\end{definition}
\section{Métrica y topología}
Sabemos que $\displaystyle \C \sim \R^{2} $ por lo que va a heredar la topología usual de $\displaystyle \R^{2} $. Recordamos que 
\[d : \C \times \C \to [0,\infty) : \left(z,w\right) \to \left|z-w\right| ,\]
es la métrica que usamos. Los abiertos de $\displaystyle \C $ estarán formados por uniones de discos abiertos:
\[D\left(z,r\right) : = \left\{ w \in \C \; : \: \left|w-z\right| < r\right\} \quad \text{y} \quad \overline{D}\left(z,r\right):= \left\{ w \in \C \; : \; \left|w-z\right|\leq r\right\}  .\]
Diremos que $\displaystyle \Omega \subset \C $ es abierto si 
\[\forall z \in \Omega, \; \exists r > 0, \; D\left(z, r\right)\subset \Omega  .\]
Diremos que la sucesión $\displaystyle \left\{ z_{n}\right\} _{n\in\N}\subset \C $ converge a $\displaystyle z_{0} \in \C $ si $\displaystyle \forall \epsilon > 0 $ existe $\displaystyle n_{0} \in \N $ tal que $\displaystyle \forall n \geq n_{0} $ se tiene que $\displaystyle \left|z_{n}-z_{0}\right|<\epsilon  $. 
En particular, dada $\displaystyle \left\{ z_{n}\right\} _{n\in\N}\subset \C $, decimos que $\displaystyle \sum^{\infty}_{n = 1}z_{n}=\infty $ si $\displaystyle \lim_{N \to \infty}\sum^{N}_{n = 1}z_{n} = z_{0} $.
\begin{definition}[Función continua]
Dada $\displaystyle f : \Omega \to \C $, diremos que $\displaystyle f $ es \textbf{continua} en $\displaystyle z_{0}\in \Omega  $ si $\displaystyle \forall \epsilon > 0 $, $\displaystyle \exists\delta > 0 $ tal que si $\displaystyle \left|z-z_{0}\right|<\delta $ entonces $\displaystyle \left|f\left(z\right)-f\left(z_{0}\right)\right|<\epsilon  $. 
\end{definition}
\begin{observation}
El espacio $\displaystyle \C $ es un espacio métrico completo. Es decir, toda sucesión de Cauchy es convergente.
\end{observation}
\begin{observation}
Las operaciones de suma y producto en $\displaystyle \C $ son polinómicas por lo que son continuas. Así, $\displaystyle \left(\C, + , \cdot \right) $ es un \textbf{cuerpo topológico}.
\end{observation}
\begin{eg}
Una consecuencia de este hecho es que si $\displaystyle z_{n} \to \alpha  $ y $\displaystyle w_{n}\to \beta  $, entonces $\displaystyle z_{n}+w_{n}\to \alpha + \beta  $. Lo mismo sucede para el producto. 
\end{eg}
\begin{theorem}[Criterio $\displaystyle M $-Weirstrass]
	Sea $\displaystyle X \neq \emptyset $, sea $\displaystyle \left(E, \left| \cdot \right|\right) $ un espacio normado completo \footnote{A estos se les llama \textbf{Espacios de Banach}.} y sea $\displaystyle f_{n} : X \to E $ una sucesión de funciones tales que existe $\displaystyle \left\{ M_{n}\right\} _{n\in\N}\subset[0,\infty) $ con $\displaystyle \left|f_{n}\left(x\right)\right|\leq M_{n} $, $\displaystyle \forall x \in X $, $\displaystyle \forall n \in \N $ y además $\displaystyle \sum^{\infty}_{n = 1}M_{n}<\infty  $. 
	Entonces existe $\displaystyle f : X \to E $ tal que $\displaystyle f\left(x\right)=\lim_{N \to \infty}\sum^{N}_{n = 1}f_{n}\left(x\right) $ uniformemente en $\displaystyle X $. Además, si $\displaystyle \left\{ f_{n}\right\} _{n\in\N} $ son continuas, $\displaystyle f $ también lo es.
\end{theorem}
\begin{proof}
Sea $\displaystyle \epsilon > 0 $. Entonces existe $\displaystyle N_{0} \in \N $ tal que si $\displaystyle k \geq N_{0} $, tenemos que $\displaystyle \sum_{n \geq k}M_{n}<\epsilon  $. Así, tenemos que $\displaystyle \forall n > m \geq N_{0} $ tenemos que 
\[ \left|\sum^{n}_{j = 1}f_{j}\left(x\right)-\sum^{m}_{j = 1}f_{j}\left(x\right)\right| \leq \sum^{n}_{j = m + 1} \left|f_{j}\left(x\right)\right| \leq \sum_{j \geq m + 1} \left|f_{j}\left(x\right)\right| \leq \sum_{j \geq m + 1}M_{j}<\epsilon .\]
Así, la sucesión $\displaystyle \left\{ \sum^{n}_{j = 1}f_{j}\left(x\right)\right\} _{n\in\N}\subset E $ es de Cauchy y por ser $\displaystyle E $ completo converge, por lo que existe
\[f\left(x\right):=\lim_{n \to \infty}\sum^{n}_{j = 1}f_{j}\left(x\right) .\]
Veamos ahora que la convergencia es uniforme. Si tomamos el límite en la inecuación anterior tenemos que 
\[ \left|f\left(x\right)-\sum^{m}_{j = 1}f_{j}\left(x\right)\right| \leq \epsilon, \; \forall x \in X  ,\]
por lo que la convergencia es uniforme.
Basta con ver que una sucesión de funciones continuas que convergen uniformemente tiene como límite una función que también es continua. Consideremos que la sucesión $\displaystyle \left\{ g_{n}\right\} _{n\in\N} $ de funciones continuas converge uniformemente a la función $\displaystyle g $. Ahora, sea $\displaystyle x_{0} \in X $ y $\displaystyle \epsilon > 0 $. Por hipótesis existe $\displaystyle n_{0}\in \N $ tal que $\displaystyle \left|g\left(x\right)-g_{n_{0}}\left(x\right)\right|<\frac{\epsilon }{3} $, $\displaystyle \forall x \in X $. 
Como $\displaystyle g_{n_{0}} $ es continua en $\displaystyle x_{0} $, tenemos que existe $\displaystyle \delta > 0 $ tal que si $\displaystyle \left|x - x_{0}\right| < \delta  $ entonces $\displaystyle \left|g_{n_{0}}\left(x\right)-g_{n_{0}}\left(x_{0}\right)\right|<\frac{\epsilon }{3} $. Así, tenemos que 
\[ \left|g\left(x\right)-g\left(x_{0}\right)\right|\leq \left|g\left(x\right)-g_{n_{0}}\left(x\right)\right| + \left|g_{n_{0}}\left(x\right)-g_{n_{0}}\left(x_{0}\right)\right|+ \left|g_{n_{0}}\left(x_{0}\right)-g\left(x_{0}\right)\right| < \frac{\epsilon }{3} + \frac{\epsilon }{3} + \frac{\epsilon }{3}=\epsilon  .\]
\end{proof}
\section{El plano complejo ampliado}
\begin{definition}[Esfera de Riemann]
	Llamaremos \textbf{esfera de Riemann} o \textbf{plano complejo ampliado} al conjunto $\displaystyle \hat{\C} := \C\cup \left\{ \infty\right\}  $ que es la compactificación de Alexandroff o de un punto del espacio topológico $\displaystyle \left(\C, \mathcal{T}_{u}\right) $. 
\end{definition}
\begin{observation}
	Es cierto que $\displaystyle \hat{\C} $ es homeomorfo a $\displaystyle \mathbb{S}^{2}= \left\{ x^{2}+y^{2}+z^{2}=1\right\}  $. 
\end{observation}
\begin{observation}
Se puede comprobar que $\displaystyle \left(x,y,z\right)\to w = \frac{x + yi}{1-z} $ es una biyección. %comprobar en casa.
Se puede comprobar que la inversa de esta aplicación está dada por

\[x = \frac{w + \overline{w}}{1 + \left|w\right|^{2}} \quad
y = \frac{w - \overline{w}}{i \left(1 + \left|w\right|^{2}\right)}\quad
z = \frac{ \left|w\right|^{2}-1}{ \left|w\right|^{2}+1} .\]
\end{observation}
\subsection{Topología de $\displaystyle \hat{\C} $}
Cualquier abierto de $\displaystyle \C $ es abierto de $\displaystyle \widehat{\C} $. Cualquier conjunto de la forma
\[ \left\{ \infty\right\} \cup \left\{ z \in \C \; : \; \left|z\right| > R\right\}  ,\]
para algún $\displaystyle R > 0 $ también es abierto.
\begin{theorem}
	La proyección esterográfica lleva las circunferencias de $\displaystyle \mathbb{S}^{2} $ ($\displaystyle \widehat{\C} $) a circunferencias de $\displaystyle \C $ o bien rectas del plano complejo (si contiene el polo norte). 
\end{theorem}
\begin{proof}
Demostrar en casa.
\end{proof}
\begin{observation}
\begin{enumerate}
	\item En $\displaystyle \hat{\C} $ decimos que $\displaystyle \left\{ z_{n}\right\} _{n\in\N}\subset \C $ converge a $\displaystyle \infty $ si $\displaystyle \left|z_{n}\right|\to \infty $.\item Con la convención $\displaystyle \frac{1}{z}=\infty $ y $\displaystyle \frac{1}{\infty}=0 $, la función $\displaystyle w = \varphi\left(z\right)=\frac{1}{z} $ es muy útil para estudiar el punto $\displaystyle \infty \in \hat{\C}$ y $\displaystyle \varphi : \hat{\C} \to \hat{\C} $ así definido es continua en $\displaystyle 0 $ y $\displaystyle \infty $ y biyectiva. 
Se puede visualizar como una rotación de $\displaystyle \pi  $ radianes en la esfera de Riemann.
Así, $\displaystyle f : \hat{\C} \to \hat{\C} $ es continua en $\displaystyle \infty $ si y solo si $\displaystyle z \to f\left(\frac{1}{z}\right) $ es continua en $\displaystyle z = 0 $.
\end{enumerate}
\end{observation}
\section{Transformaciones de Möbius}
\begin{definition}[Transformación de Möebius]
	Las funciones de la forma $\displaystyle f\left(z\right)=\frac{az + b}{cz + d} $ con $\displaystyle a,b,c,d \in \C $ tales que $\displaystyle \begin{vmatrix} a & b \\ c & d \end{vmatrix} \neq 0 $, se llaman \textbf{transformaciones de Möbius}. 
\end{definition}
\begin{notation}
	Escribiremos, para $\displaystyle A = \begin{pmatrix} a & b \\ c & d \end{pmatrix} $ con $\displaystyle \det A \neq 0 $, $\displaystyle f\left(z\right)=\frac{az+b}{cz+d}=:f_{A}\left(z\right) $.
\end{notation}
\begin{observation}
La hipótesis de que $\displaystyle \det A \neq 0 $ implica que $\displaystyle f $ no es constante.
\end{observation}
\begin{eg}
\begin{enumerate}
\item Si ponemos $\displaystyle c = 0 $, $\displaystyle d = 1 $ y $\displaystyle a \neq 0 $, tenemos que $\displaystyle f\left(z\right)=az + b $ es una transformación afín. En particular, traslaciones y homotecias son transformaciones de Möbius. 
\item $\displaystyle f\left(z\right)=\frac{1}{z} $ las llamamos inversiones.
\end{enumerate}
\end{eg}
\begin{prop}
\begin{enumerate}
\item Las transformaciones de Möbius $\displaystyle f_{A} $ se pueden extender a $\displaystyle \hat{f}_{A} : \hat{\C} \to \hat{\C} $. Diremos que $\displaystyle \hat{f}_{A}\left(\infty\right) $ si $\displaystyle f_{A} $ es afín y si $\displaystyle c \neq 0 $ ponemos $\displaystyle \hat{f}\left(\infty\right)=\frac{a}{c} $, logrando así que $\displaystyle \hat{f}_{A} $ sea continua en $\displaystyle \infty \in \hat{\C}$. 
	Asimismo, definimos $\displaystyle \hat{f}_{A}\left(-\frac{d}{c}\right)=\infty \in \hat{\C} $. Así, $\displaystyle \hat{f}_{A} $ es continua.
\item $\displaystyle f_{A} $ es biyectiva y su inversa es $\displaystyle f_{A^{-1}} $.
\item $\displaystyle f_{A}\circ f_{B} = f_{A \cdot B} $. 
\end{enumerate}
\end{prop}
\begin{proof}
\textbf{(1)} se sigue de \textbf{(2)} ya que $\displaystyle \GL_{2}\left(\C\right) $ es un grupo. 
\end{proof}

\begin{observation}
El conjunto de transformaciones de Möbius forma un grupo bajo composición. La propiedad de grupo se hereda del hecho que $\displaystyle \GL_{2}\left(\C\right) $ es un grupo. 
\end{observation}
\begin{observation}
Consideremos $\displaystyle f_{A}\left(z\right)=\frac{az + b}{cz + d} $. Dado $\displaystyle \lambda \in \C^{\times} $, tenemos que 
\[f_{\lambda A}\left(z\right)= f_{A}\left(z\right) .\]
Así, realmente el grupo que actua es $\displaystyle \mathbb{P}\GL_{2}\left(\C\right)= \GL_{2}\left(\C\right)/\C $. 
\end{observation}
\begin{theorem}
Dados $\displaystyle z_{1}, z_{2}, z_{3} \in \hat{\C} $ distintos y otros $\displaystyle w_{1}, w_{2}, w_{3} \in \hat{\C} $ distintos, existe una única $\displaystyle \hat{f}_{A} : \hat{\C} \to \hat{\C} $ tal que $\displaystyle f_{A}\left(z_{j}\right)=w_{j} $ con $\displaystyle j = 1, 2, 3 $.
\end{theorem}
\begin{proof}
Nos basta con encontrar $\displaystyle f_{A} $ tal que 
\[
\begin{cases}
z_{1} \to 0 \\
z_{2} \to 1 \\
z_{3} \to \infty
\end{cases}
.\]
En efecto, si podemos construir $\displaystyle f_{A} $ también podemos construir $\displaystyle f_{B} $ tal que 
\[
\begin{cases}
w_{1} \to 0 \\
w_{2} \to 1 \\
w_{3} \to \infty
\end{cases}
.\]
Así, tomando $\displaystyle f_{B^{-1}A} $ hemos encontrado la función que buscábamos. Distinguirmos ahora varios casos:
\begin{description}
	\item[Caso 1.] Si $\displaystyle \infty \not\in \left\{ z_{1}, z_{2}, z_{3}\right\}  $, podemos coger
		\[f\left(z\right):= \frac{z - z_{1}}{z-z_{3}} \cdot \frac{z_{2}-z_{3}}{z_{2}-z_{1}} .\]
	\item[Caso 2.] Si $\displaystyle \infty \in \left\{ z_{1}, z_{2}, z_{3}\right\}  $, supongamos que $\displaystyle z_{1}= \infty $. Ponemos pues
		\[g_{1}\left(z\right)=\frac{z_{2}-z_{3}}{z-z_{3}} .\]
		Si $\displaystyle z_{2}=\infty $ cogemos 
		\[g_{2}\left(z\right)=\frac{z-z_{1}}{z-z_{3}} .\]
	Finalmente, si $\displaystyle z _{3}= \infty $, tendremos que 
	\[g_{3}\left(z\right)=\frac{z-z_{1}}{z_{2}-z_{1}} .\]
\end{description}
Veamos ahora la unicidad. Supongamos que existen dos soluciones $\displaystyle f $ y $\displaystyle g $. Necesitamos ver primero que si 
\[h : 
\begin{cases}
0 \to 0 \\
1 \to 1 \\
\infty \to \infty
\end{cases}
,\]
entonces $\displaystyle h\left(z\right)= z $. Esto es muy fácil de comprobar. Sea $\displaystyle \varphi\circ f^{-1}\circ g \circ \varphi^{-1} $ tal que 
\[\varphi:
\begin{cases}
z_{1} \to 0 \\
z_{2} \to 1 \\
z_{3} \to \infty
\end{cases}
.\]
Vemos que 
\[\varphi \circ f^{-1}\circ g\circ \varphi^{-1} : 
\begin{cases}
0 \to 0 \\
1 \to 1 \\
\infty \to \infty 
\end{cases}
.\]
Obligatoriamente debe ser que $\displaystyle \varphi\circ f^{-1}\circ g\circ\varphi^{-1}=id\circ \varphi $, por lo que $\displaystyle f^{-1}\circ g = id $ y tenemos que $\displaystyle f = g $.
\end{proof}
\begin{theorem}
Toda función de Möbius $\displaystyle f_{A} $ es composición de homotecias, traslaciones e inversiones. 
\end{theorem}
\begin{proof}
Distinguimos dos casos:
\begin{description}
\item[Caso 1.] Si $\displaystyle f\left(\infty\right)=\infty $ entonces $\displaystyle f $ es afín. 
\item[Caso 2.] Supongamos que $\displaystyle f\left(\infty\right)\neq \infty  $, entonces debe ser que $\displaystyle f\left(z\right)=\frac{az + b}{cz + d} $ con $\displaystyle c \neq 0 $. Sin pérdida de generalidad podemos asumir que $\displaystyle c= 1 $ y obtenemos que $\displaystyle f\left(z\right)=\frac{az + b}{z + d}=a + \frac{b-ad}{z + d} $. 
\end{description}
\end{proof}
\begin{theorem}
Toda transformación de Möbius $\displaystyle \hat{f}_{A}: \hat{\C} \to \hat{\C} $ lleva circunferencias de $\displaystyle \hat{\C} $ a circunferencias de $\displaystyle \hat{\C} $. 
\end{theorem}
\begin{proof}
Aplicando el teorema anterior nos dice que basta con comprobar que el resultado es cierto para traslaciones, homotecias e inversiones.
% Ejercicio de la hoja 2
\end{proof}

